\section{开放问题：下一代 Agent 安全研究的技术缺口}

课程虽然给了清晰框架，但仍留下大量未解问题。下面从研究与工程交叉视角总结最关键的技术缺口。

\subsection{缺口一：可证明安全与高灵活性的张力}
Agent 价值来自灵活性，但可证明安全通常需要收窄行为空间。如何在开放任务中同时保持高 utility 与可验证安全，是目前最大的基础矛盾。现有方法往往二选一：要么很安全但任务能力下降，要么能力强但缺少形式化保证。

\begin{importantbox}{研究方向}
将 ``高危动作'' 抽象为可证明约束，把 ``低危动作'' 保留给模型自由探索。也就是说，不追求对全行为空间证明安全，而追求对关键后果路径证明安全。
\end{importantbox}

\subsection{缺口二：跨环境泛化与跨任务泛化}
很多系统在单环境 benchmark 上表现良好，但换工具、换数据分布、换任务模板后安全性能快速退化。这说明当前防御很多是 ``环境特化防御''，而非真正的泛化防御。

\begin{knowledgebox}{可操作评测建议}
在评测中显式加入：
\begin{itemize}
    \item \textbf{Cross-environment split}：训练和测试环境隔离。
    \item \textbf{Cross-tool split}：训练工具集合与测试工具集合不重叠。
    \item \textbf{Cross-policy split}：策略模板变化后重新评估鲁棒性。
\end{itemize}
\end{knowledgebox}

\subsection{缺口三：安全基准的标准协议仍不统一}
尽管课程提出 AAA 思路，但行业还缺统一协议与统一数据格式。没有协议，生态难互通；没有互通，防御研究难积累。安全 benchmark 的 ``可持续运营'' 也是难点，需要处理泄题、版本漂移、指标膨胀和社区治理。

\subsection{缺口四：攻防不对称下的资源配置}
讲者在网络安全章节指出，短期内 AI 可能先帮助攻击者。对防守团队而言，核心问题不是 ``能不能做最强防御''，而是 ``在有限算力、人力、响应时间下，防哪里最划算''。这需要把安全研究和安全经济学结合。

\begin{warningbox}{容易被忽视的现实约束}
如果防御机制引入过高时延或过高误报，业务方会绕开安全机制。最终系统会回到 ``名义安全、实际裸奔''。因此安全方案必须同时优化风险降低与业务可接受性。
\end{warningbox}

\subsection{缺口五：人机协同边界如何设计}
Human-in-the-loop 不是简单加一个确认弹窗。真正的问题是：什么动作必须人工确认、人工在多大负载下仍能可靠判断、如何避免 ``确认疲劳''。这需要行为科学、界面设计和系统安全共同参与。

\begin{table}[H]
\centering
\begin{tabular}{p{2.9cm}p{3.6cm}p{3.9cm}}
\toprule
\textbf{开放问题} & \textbf{当前瓶颈} & \textbf{潜在突破方向} \\
\midrule
可证明安全 & 行为空间过大难以全量证明 & 关键路径可证明 + 分层约束 \\
跨环境鲁棒性 & 防御过度依赖特定环境 & 跨环境训练与评测协议 \\
标准化评测 & 协议与指标碎片化 & AAA/MCP 风格统一接口 \\
攻防资源分配 & 防守成本高、优先级不清晰 & 风险经济模型 + 动态预算 \\
人机协同 & 审批疲劳与误判 & 风险分级审批与自适应交互 \\
\bottomrule
\end{tabular}
\caption{Agent 安全下一阶段的核心技术缺口}
\end{table}

\subsection{本章小结}
Agent 安全的下一步，不是继续堆单点 patch，而是解决 ``可验证性、泛化性、标准化、经济性、人机协同'' 五个结构性问题。这些问题决定了 Agent 能否在高风险场景长期可用。

